\documentclass[conference]{IEEEtran}
\IEEEoverridecommandlockouts

\usepackage{cite}
\usepackage{amsmath,amssymb,amsfonts}
\usepackage{graphicx}
\usepackage{textcomp}
\usepackage{xcolor}
\usepackage{booktabs}
\usepackage{url}
\usepackage{hyperref}
\usepackage{multirow}
\usepackage{array}

\def\BibTeX{{\rm B\kern-.05em{\sc i\kern-.025em b}\kern-.08em
    T\kern-.1667em\lower.7ex\hbox{E}\kern-.125emX}}

\begin{document}

\title{CircuitBench: Benchmarking Vision-Language Models\\
for Electronic Circuit Reasoning and Fault Diagnosis}

\author{\IEEEauthorblockN{Sasha Sam}
\IEEEauthorblockA{\textit{Purdue University}\\
West Lafayette, Indiana, USA}}

\maketitle

% ===========================================================================
\begin{abstract}
Vision-language models (VLMs) are increasingly proposed for use in
engineering education, including as conversational laboratory assistants
capable of interpreting circuit diagrams and diagnosing faults.
Before such systems are deployed, their reliability on structured
electrical-engineering tasks must be characterized.
We present \emph{CircuitBench}, a source-backed benchmark for evaluating
VLMs on electronic-circuit understanding, reasoning, and fault diagnosis,
derived from two publicly available datasets: the AITEE Electrical Circuit
Dataset (828 annotated schematic images~\cite{knievel2025aitee_dataset})
and the CGHD hand-drawn circuit dataset~\cite{thoma2025cghd}.
The benchmark comprises 8,598 items across three task categories:
component recognition (Task~A), connectivity and topology reasoning
(Task~B), and fault diagnosis (Task~E), with a 2,581-item held-out
test split.
We evaluate three locally-hosted open-weight VLMs under a standardized
zero-shot prompting protocol on a development-split sample.
The best-performing model, Gemma~3~4B, achieves 38.3\% overall accuracy
($n{=}81$), with markedly higher performance on binary yes/no subtasks
(fault presence: 87.5\%; topology structure: 87.5\%) than on counting or
free-text tasks (node count: 0\%; circuit description: 0\%).
LLaVA~7B achieves 23.3\% and LLaVA-Phi3 achieves 8.0\%.
All models fail completely on connectivity-counting and free-text
description tasks.
These results suggest that current small open-weight VLMs are not
sufficiently reliable for autonomous use in electronics education or
circuit debugging, and that fault detection and topology reasoning
present distinct capability profiles worth investigating further.
\end{abstract}

\begin{IEEEkeywords}
vision-language models, multimodal reasoning, electrical circuits,
fault diagnosis, engineering education, benchmark
\end{IEEEkeywords}


% ===========================================================================
\section{Introduction}
\label{sec:intro}

Vision-language models have demonstrated strong performance on visual
question answering, scientific diagram interpretation, and chart
understanding~\cite{liu2023mmbench,lu2022scienceqa,masry2022chartqa}.
As these models improve, there is growing interest in applying them to
engineering education---for example, as conversational laboratory
assistants that can identify circuit components, trace signal paths, and
diagnose wiring errors~\cite{knievel2025aitee_paper}.

However, circuits are not generic images.
Interpreting an electrical schematic requires recognizing specific symbols
(resistors, voltage sources, grounds), understanding how those symbols are
connected by wires and nodes, reasoning about series and parallel
topologies, and---for fault diagnosis---identifying which connections or
components violate expected circuit structure.
An incorrect diagnosis from an AI assistant could mislead a student
attempting to debug a real circuit.

Before such systems are deployed, it is important to characterize
whether contemporary VLMs can perform these tasks reliably.
Prior benchmarks focus on natural images~\cite{liu2023mmbench}, scientific
figures~\cite{lu2022scienceqa}, and charts~\cite{masry2022chartqa}, but
systematic evaluation on electrical circuit diagrams has received limited
attention.
The AITEE project introduced annotated educational circuit images and an
agentic tutoring pipeline~\cite{knievel2025aitee_paper,knievel2025aitee_dataset}
but did not conduct a cross-model VLM benchmark.
The CGHD dataset provides rich annotations for hand-drawn
circuits~\cite{thoma2021cghd_paper,thoma2025cghd} but is a recognition
resource, not an evaluation benchmark.

We introduce \emph{CircuitBench}, a source-backed benchmark that fills
this gap.
It spans three circuit-reasoning task categories, is derived from
two publicly licensed datasets, and includes controlled fault-injection
items with deterministic ground truth.
We evaluate three open-weight VLMs under identical zero-shot conditions
and analyze their failure modes.

\textbf{Contributions:}
\begin{itemize}
  \item CircuitBench, a benchmark of 8,598 items derived from AITEE and
        CGHD with full provenance metadata and a frozen 2,581-item test split;
  \item a controlled fault-injection protocol producing 798 fault-injected
        images with deterministic labels from 200 source circuits;
  \item an empirical comparison of three open-weight VLMs under
        standardized zero-shot prompting; and
  \item an error analysis identifying recurring failure modes across models
        and task categories.
\end{itemize}


% ===========================================================================
\section{Related Work}
\label{sec:related}

\subsection{Vision-Language Models and Multimodal Reasoning}

Large vision-language models such as GPT-4o~\cite{openai2024gpt4o},
Gemini~\cite{google2024gemini}, and LLaVA~\cite{liu2023llava}
combine a visual encoder with a language model decoder trained on
image-text pairs.
MMBench~\cite{liu2023mmbench} and ScienceQA~\cite{lu2022scienceqa}
provide broad multimodal evaluation, while ChartQA~\cite{masry2022chartqa}
focuses on chart comprehension.
Structured technical diagrams such as circuit schematics differ
substantially from the natural images and charts that dominate these
benchmarks, and their evaluation remains underexplored.

\subsection{AI for Engineering Education}

AI tutoring systems have been applied across STEM subjects, with
vision-language models enabling diagram-based question
answering~\cite{lu2022scienceqa}.
The AITEE system~\cite{knievel2025aitee_paper} is the most directly
relevant recent work: it introduces 828 annotated first-semester
circuit images and proposes an agentic tutoring pipeline using
image-to-netlist conversion.
CircuitBench uses AITEE as its primary data source.

\subsection{Circuit Diagram Recognition and Understanding}

Automated circuit interpretation has been studied as a pipeline:
detect symbols, extract connectivity, generate a netlist.
Thoma et al.~\cite{thoma2021cghd_paper} introduced the CGHD dataset
and demonstrated deep-learning-based component localization in hand-drawn
circuit images.
Alakhras et al.~\cite{alakhras2026sina} showed that combining deep
learning, connected-component labeling, and OCR achieves 96.67\%
netlist generation accuracy on digital schematics.
Nguyen et al.~\cite{nguyen2024circuitparse} applied Graph Attention
Networks to netlist parsing of integrated-circuit images.
These works extract structured representations from circuit images;
CircuitBench instead evaluates whether VLMs can answer natural-language
reasoning questions \emph{without} access to extracted annotations.

\subsection{Technical Visual-Reasoning Benchmarks}

ScienceQA~\cite{lu2022scienceqa} and ChartQA~\cite{masry2022chartqa}
evaluate multimodal reasoning on science figures and charts,
respectively, but neither includes electrical schematics or circuit-level
reasoning tasks.
To our knowledge, CircuitBench is the first benchmark specifically
designed to evaluate VLMs on component recognition, topology reasoning,
and fault diagnosis from electrical circuit diagrams.


% ===========================================================================
\section{CircuitBench}
\label{sec:benchmark}

\subsection{Source Datasets}

\textbf{AITEE.}
The AITEE Electrical Circuit Dataset~\cite{knievel2025aitee_dataset}
contains 828 annotated schematic images of first-semester electrical
engineering circuits, including series, parallel, mixed series-parallel,
voltage-divider, Wheatstone bridge, delta, and star topologies.
Images are accompanied by YOLO-format annotations covering 11 classes
(edge/junction, ground, current source, node point, resistor, voltage
source, and identifier classes).
The dataset is released under the Apache~2.0 license.

\textbf{CGHD.}
The CGHD dataset (Version~16, DOI:~\texttt{10.5281/zenodo.17469897})~%
\cite{thoma2025cghd} provides 3,269 annotated hand-drawn circuit
diagram images contributed by 31 main drafters, plus additional
real-world examination images.
It includes 248,020 bounding-box annotations across 59 object classes.
CGHD Version~16 is released under CC~BY~4.0.

\subsection{Data Selection}

AITEE images were included if the corresponding label file contained at
least one resistor annotation, ensuring non-trivial circuit content.
All 828 available images met this criterion.
For fault injection, images with at least three resistors and three
edge/junction annotations were selected to ensure structural
meaningfulness after modification; this yielded 200 source circuits.

\subsection{Task Taxonomy}

CircuitBench evaluates three task categories:

\textbf{Task A --- Component recognition} covers three subtasks:
\emph{component\_presence} (binary: is a given component type present?),
\emph{component\_count} (integer: how many of a given type?), and
\emph{component\_enumeration} (free-text: list all types with counts).

\textbf{Task B --- Connectivity and topology reasoning} covers three
subtasks: \emph{topology\_structure} (binary: are two components in
series or parallel?), \emph{node\_count} (integer: how many circuit
nodes?), and \emph{component\_total} (integer: total component count).

\textbf{Task E --- Fault diagnosis} covers three subtasks:
\emph{fault\_presence} (binary: is a fault present?),
\emph{fault\_diagnosis} (free-text: identify and describe the fault),
and \emph{circuit\_description} (free-text: describe all component types
and general topology).

\subsection{Controlled Fault Generation}

Task~E uses modified circuit images in which a single fault has been
deterministically introduced by one of two operations:

\textbf{Open circuit:} a white rectangle is drawn over the wire segment
between two nearby edge/junction nodes, creating a broken connection.

\textbf{Disconnected component:} a randomly selected resistor's bounding
box (plus a 5-pixel margin) is filled white, removing the component.

Only modified copies are used; source images are never altered.
Each fault produces both a presence (E1) and a diagnosis (E2) item.
All transformation details are recorded in the benchmark manifest.

\subsection{Ground-Truth Construction}

Binary and integer ground truth for Tasks~A and~B is derived
deterministically from AITEE's YOLO annotations.
Fault ground truth is specified by construction: the fault type is known
exactly from the transformation applied.
Free-text responses are scored by keyword matching, described in
Section~\ref{sec:setup}.

\subsection{Benchmark Statistics and Splits}

Table~\ref{tab:dataset} summarizes the benchmark.
Items are split 70/30 test/dev by source image, keeping all items
from the same source circuit in the same split to prevent leakage.

\begin{table}[t]
\centering
\caption{CircuitBench Dataset Composition}
\label{tab:dataset}
\begin{tabular}{llrr}
\toprule
\textbf{Task} & \textbf{Subtask} & \textbf{Dev} & \textbf{Test} \\
\midrule
\multirow{3}{*}{A --- Component}
  & component\_presence    & 1{,}158 &   498 \\
  & component\_count       &    975  &   410 \\
  & component\_enumeration &    579  &   249 \\
  & \textit{A total}       & \textit{2{,}712} & \textit{1{,}157} \\
\midrule
\multirow{3}{*}{B --- Topology}
  & topology\_structure    &    577  &   249 \\
  & node\_count            &    577  &   249 \\
  & component\_total       &    579  &   249 \\
  & \textit{B total}       & \textit{1{,}733} & \textit{747} \\
\midrule
\multirow{3}{*}{E --- Fault}
  & fault\_presence        &    859  &   368 \\
  & circuit\_description   &    433  &   190 \\
  & fault\_diagnosis       &    280  &   119 \\
  & \textit{E total}       & \textit{1{,}572} & \textit{677} \\
\midrule
\textbf{Total} & & \textbf{6{,}017} & \textbf{2{,}581} \\
\bottomrule
\end{tabular}
\vspace{2pt}

\footnotesize{Task~E includes 798 fault-injected items from 200 source
circuits (up to two fault types each, yielding E1 and E2 items per fault).}
\end{table}

\subsection{Licensing and Attribution}

AITEE images are used under the Apache~2.0 license.
CGHD images are used under CC~BY~4.0.
CircuitBench benchmark items (questions, ground truth, fault-injected
images) are original contributions of this work.
The AITEE and CGHD dataset creators are not authors of CircuitBench.


% ===========================================================================
\section{Experimental Setup}
\label{sec:setup}

\subsection{Evaluated Models}

Table~\ref{tab:models} lists the three models evaluated.
All are open-weight models run locally via Ollama using 4-bit quantized
GGUF weights, requiring no API keys.

\begin{table}[t]
\centering
\caption{Evaluated Models}
\label{tab:models}
\begin{tabular}{lll}
\toprule
\textbf{Model} & \textbf{Parameters} & \textbf{Quant.} \\
\midrule
Gemma 3 4B~\cite{gemmateam2025gemma3}  & 4.3B & Q4\_K\_M \\
LLaVA 7B~\cite{liu2023llava}           & 7B   & Q4\_K\_M \\
LLaVA-Phi3~\cite{liu2023llava,abdin2024phi3} & 4B & Q4\_K\_M \\
\bottomrule
\end{tabular}
\vspace{2pt}

\footnotesize{All models accessed via Ollama on an Intel Core Ultra CPU
with 16\,GB RAM. Model sizes: Gemma~3~4B (3.3\,GB), LLaVA~7B (4.7\,GB),
LLaVA-Phi3 (2.9\,GB).}
\end{table}

\subsection{Prompting Protocol}

Each item is presented as a single user turn: the circuit image followed
by a task-specific prompt template.
Templates are identical across all models; no system prompt is used.
Example templates:

\textbf{Fault presence (E1):}
\textit{``Analyze the electrical circuit shown in the image.
Is there a circuit fault present? Answer Yes or No.''}

\textbf{Topology structure (B2):}
\textit{``Examine the electrical circuit diagram carefully.
Pay attention to how the components are connected.
\{question\}
Respond with only `Yes' or `No'. Do not include any explanation.''}

\textbf{Component count (A2):}
\textit{``Examine the electrical circuit diagram in the provided image.
\{question\}
Respond with only the number. Do not include any explanation.''}

All prompt templates are included in the benchmark release.

\subsection{Inference Parameters}

Temperature: 0; maximum output tokens: 256; no in-context examples
(zero-shot).
A 0.5-second delay is applied between requests.

\subsection{Evaluation Metrics}

\textbf{Binary tasks} (E1, B2, A1): accuracy based on whether the
first token of the response matches the ground-truth ``yes'' or ``no''
(case-insensitive).

\textbf{Integer tasks} (A2, B1): accuracy based on whether the first
integer token in the response matches the ground-truth count.

\textbf{Free-text tasks} (E2, A3, E3): keyword matching over component
names and counts from the ground truth.
A response is marked correct only if all required keywords are present.
This scorer is conservative and may undercount genuinely correct
responses phrased differently.

\subsection{Evaluation Sample}

The evaluation presented here uses a development-split sample:
81 items for Gemma~3~4B, 50 items for LLaVA-Phi3, and 31 items for
LLaVA~7B, drawn from the 6,017-item development split.
All items were sampled before model evaluation; no test-split items
were used.
The 2,581-item frozen test split is reserved for future full-scale
evaluation.


% ===========================================================================
\section{Results}
\label{sec:results}

\subsection{Overall Performance}

Table~\ref{tab:overall} reports overall and per-task accuracy for all
three models.
Gemma~3~4B achieves the highest overall accuracy at 38.3\%,
followed by LLaVA~7B at 23.3\%, and LLaVA-Phi3 at 8.0\%.
No model exceeds 40\% on any task category, and all models perform
at or near 0\% on Task~B (topology/connectivity reasoning) in
aggregate---largely driven by complete failure on node-count and
component-total subtasks.

\begin{table}[t]
\centering
\caption{Overall and Per-Task Accuracy (Development Split)}
\label{tab:overall}
\begin{tabular}{lcccc}
\toprule
\textbf{Model} & \textbf{Overall} & \textbf{Task A} & \textbf{Task B} & \textbf{Task E} \\
\midrule
Gemma 3 4B  & 38.3\% & 38.6\% & 29.2\% & 53.8\% \\
LLaVA 7B    & 23.3\% & 26.7\% & 33.3\% &  0.0\% \\
LLaVA-Phi3  &  8.0\% & 15.4\% &  0.0\% &  0.0\% \\
\midrule
\textit{n (items)} & & \textit{81 / 30 / 50} & \textit{81 / 30 / 50} & \textit{81 / 30 / 50} \\
\bottomrule
\end{tabular}
\vspace{2pt}
\footnotesize{Model order: Gemma~3~4B / LLaVA~7B / LLaVA-Phi3.}
\end{table}

\subsection{Performance by Subcategory}

Table~\ref{tab:subtask} provides a detailed breakdown.
Several patterns are immediately apparent.

\begin{table}[t]
\centering
\caption{Accuracy by Task Subcategory}
\label{tab:subtask}
\begin{tabular}{lccc}
\toprule
\textbf{Subcategory} & \textbf{Gemma 3 4B} & \textbf{LLaVA 7B} & \textbf{LLaVA-Phi3} \\
\midrule
\multicolumn{4}{l}{\textit{Task A --- Component Recognition}} \\
component\_presence    & 47.1\% & 16.7\% & 40.0\% \\
component\_count       & 26.3\% & 50.0\% &  0.0\% \\
component\_enumeration & 50.0\% &  0.0\% &  0.0\% \\
\midrule
\multicolumn{4}{l}{\textit{Task B --- Topology Reasoning}} \\
topology\_structure    & 87.5\% &100.0\% &  0.0\% \\
node\_count            &  0.0\% &  0.0\% &  0.0\% \\
component\_total       &  0.0\% &  0.0\% &  0.0\% \\
\midrule
\multicolumn{4}{l}{\textit{Task E --- Fault Diagnosis}} \\
fault\_presence        & 87.5\% &  0.0\% &  0.0\% \\
fault\_diagnosis       &  0.0\% &  0.0\% &  0.0\% \\
circuit\_description   &  0.0\% &  0.0\% &  0.0\% \\
\bottomrule
\end{tabular}
\end{table}

\textbf{Binary yes/no tasks are disproportionately easy.}
Both Gemma~3~4B and LLaVA~7B achieve 87.5\% and 100\% respectively
on topology\_structure, and Gemma~3~4B achieves 87.5\% on
fault\_presence.
These are binary forced-choice tasks; performance may partly reflect
response-format bias rather than genuine circuit understanding.

\textbf{Counting tasks are inconsistent.}
LLaVA~7B achieves 50\% on component\_count while Gemma~3~4B achieves
26.3\% and LLaVA-Phi3 0\%.
However, all three models score 0\% on node\_count and
component\_total, suggesting that counting nodes or total components
from circuit images is beyond current local VLM capability.

\textbf{Free-text tasks fail universally.}
All three models score 0\% on fault\_diagnosis and circuit\_description.
This is partly a scorer artifact (strict keyword matching), but manual
inspection confirms that models generally either hallucinate component
types, describe non-existent circuit features, or produce responses
that do not match the expected component-name-plus-count format.

\textbf{Fault diagnosis (free-text) vs.\ fault presence (binary).}
The contrast between Gemma~3~4B's 87.5\% on fault\_presence and
0\% on fault\_diagnosis is striking.
The model can detect that \emph{something} is wrong but cannot describe
\emph{what} is wrong in a scorable format.


% ===========================================================================
\section{Error Analysis}
\label{sec:error}

Manual inspection of a sample of incorrect responses reveals the
following recurring patterns.

\textbf{Format non-compliance.}
LLaVA-Phi3 frequently ignores instructions to answer with only
``Yes'', ``No'', or a number, instead producing verbose descriptions
beginning with phrases like ``The circuit shows\ldots''
This explains its 0\% on topology\_structure despite LLaVA~7B (same
architecture, larger model) achieving 100\%---the smaller Phi3 base
model is less instruction-following.

\textbf{Counting errors.}
For component\_count, models typically answer with a value 1--3 away
from the correct count.
Gemma~3~4B tends to undercount (e.g., reporting 2 resistors when 4
are present), consistent with failing to distinguish closely placed
or overlapping symbols.
LLaVA~7B is more accurate on this subtask (50\%) but inconsistent
across circuit types.

\textbf{Node-count failure.}
All models score 0\% on node\_count.
Node counting requires understanding circuit topology globally---tracing
wires to identify all distinct electrical nodes---which appears to
exceed the spatial reasoning capability of all evaluated models in
this zero-shot setting.

\textbf{Hallucinated fault descriptions.}
On fault\_diagnosis items, Gemma~3~4B typically produces responses
describing a circuit fault in general terms (e.g., ``open circuit
detected'') without correctly identifying the specific affected
component or location.
This is semantically close to correct but fails strict keyword scoring.
A rubric-based human evaluation would likely credit partial understanding
for these responses.

\textbf{Component enumeration format mismatch.}
On circuit\_description, all models produce syntactically correct
descriptions (e.g., ``The circuit contains three resistors and one
voltage source'') but fail the scorer because the expected format
is ``Component~(N)'' notation.
These responses represent a scorer limitation as much as a model
failure and should be graded manually before drawing conclusions about
this subcategory.


% ===========================================================================
\section{Discussion}
\label{sec:discussion}

The results reveal a clear capability hierarchy among the three
evaluated models (Gemma~3~4B $>$ LLaVA~7B $>$ LLaVA-Phi3) and a
consistent pattern across all models: binary classification tasks are
handled far better than counting or free-text generation tasks.

The gap between fault\_presence (87.5\% for Gemma~3~4B) and
fault\_diagnosis (0\%) is the most practically significant finding.
A system that detects a fault but cannot explain it provides limited
value to a student debugging a circuit.
This suggests that current small open-weight VLMs may be usable as a
rough ``something looks wrong'' detector but not as a reliable
diagnostic tool.

The universal failure on node\_count (0\% across all models) points to
a fundamental gap in circuit-graph understanding.
Counting nodes requires tracing wire connectivity across the image---a
form of spatial relational reasoning that appears absent in all three
models tested.
This is particularly relevant for AI tutoring applications where
Kirchhoff's current law analysis requires correct node identification.

LLaVA-Phi3's near-complete failure (8\% overall, 0\% on all tasks
except component\_presence) appears primarily attributable to
instruction-following breakdown: the 4B Phi3 base model frequently
ignores format constraints, producing verbose answers where single
tokens are expected.
This is a practical deployment concern independent of circuit
understanding.

The benchmark's current evaluation uses development-split samples
of 30--81 items per model, which is sufficient to identify broad
patterns but insufficient for fine-grained statistical comparisons
between models.
The frozen 2,581-item test split is reserved for a future evaluation
that includes cloud-API models (GPT-4o, Gemini, Claude) for comparison
against the open-weight baselines established here.


% ===========================================================================
\section{Limitations}
\label{sec:limitations}

\begin{enumerate}

\item \textbf{Small evaluation sample.}
Development-split samples of 30--81 items per model support pattern
identification but not rigorous statistical model comparison.
Full test-split evaluation is needed for final conclusions.

\item \textbf{Open-weight models only.}
The evaluated models are small, locally-hosted, 4-bit quantized.
State-of-the-art cloud-API VLMs (GPT-4o, Gemini 2.0 Flash,
Claude Sonnet) may perform substantially better.

\item \textbf{Conservative free-text scorer.}
Keyword matching undercounts partially correct free-text responses.
Subcategories including fault\_diagnosis and circuit\_description
require human review to calibrate.

\item \textbf{Benchmark source.}
CircuitBench is derived from AITEE and CGHD, which were not designed
as VLM evaluation resources.
The circuit distribution reflects what is available in these datasets.

\item \textbf{Circuit complexity.}
AITEE depicts first-semester educational circuits.
Results may not generalize to more complex topologies.

\item \textbf{Model contamination.}
AITEE and CGHD are publicly available.
Contemporary VLMs may have encountered their images during pretraining,
and high accuracy should not be interpreted as pure unseen
generalization.

\item \textbf{Synthetic faults.}
The two fault types (open circuit, disconnected component) are
controlled synthetic modifications.
Real student wiring mistakes may exhibit different visual
characteristics.

\item \textbf{Zero-shot prompting only.}
Performance may improve with few-shot examples or chain-of-thought
prompting; this is not explored here.

\end{enumerate}


% ===========================================================================
\section{Conclusion}
\label{sec:conclusion}

We introduced CircuitBench, a source-backed benchmark of 8,598 items
for evaluating VLMs on electronic-circuit understanding, reasoning, and
fault diagnosis, derived from the AITEE and CGHD datasets.
Evaluation of three locally-hosted open-weight VLMs reveals a
consistent pattern: binary yes/no tasks (fault presence, topology
structure) are handled reasonably well by the best model (Gemma~3~4B:
87.5\%), while counting tasks, free-text diagnosis, and node
identification fail near-universally across all models.
Overall accuracy ranges from 8.0\% (LLaVA-Phi3) to 38.3\%
(Gemma~3~4B)---well below the reliability threshold required for
autonomous instructional use.

These results suggest that small open-weight VLMs in their current
form are not suitable as autonomous circuit-debugging assistants,
though they may be useful as coarse anomaly detectors or
component-presence checkers in human-in-the-loop workflows.

Future work includes full test-split evaluation with cloud-API models,
extending the benchmark to hand-drawn circuit representations (Task~F),
adding quantitative reasoning tasks (Task~D), and calibrating the
free-text scorer against human judgments.


% ===========================================================================
\section*{Data and Code Availability}

The CircuitBench benchmark manifest (JSONL), evaluation scripts,
scored results, and fault-injected images are available at
\url{https://github.com/sasha-sam/circuitbench}.
Source datasets: AITEE~\cite{knievel2025aitee_dataset} and
CGHD~\cite{thoma2025cghd}.


% ===========================================================================
\section*{AI Assistance Disclosure}

Generative AI tools, including Claude (Anthropic), were used to assist
with code development, manuscript organization, and language refinement
during preparation of this work.
All research questions, experimental design choices, data-source
selection, source verification, model evaluations, result interpretation,
citation verification, and final manuscript content were reviewed and
validated by the author.
AI systems were not treated as authors.


% ===========================================================================
\bibliographystyle{IEEEtran}
\bibliography{references}

\end{document}
